Seven Consecutive Comoving Multiples of the Tan-Arctan-τ Angular Separators: From Triangular Symmetry to Residual Six-Sliced Hexagon and the Propagation of \(E_8\) Petrie Rings

The continuous tan-arctan-τ family supplies a microscopic angular separator whose seventh multiple, when reduced modulo the full turn \(2\pi\), lies sufficiently near the ideal hexagonal generator to propagate a self-consistent, comoving radial-fold symmetry. Once that radial fold is completed, the residual six-sliced hexagon it carries continues to act on the circle as a dense, non-periodic orbit; the same orbit, further densified, generates the concentric Petrie rings that approximate the outer \(E_8\) shell of the nested hypercomplex projection.

Begin with the continuous source. The Unique Real Tau Radian \(\operatorname{Re}(\tau)=\arctan(2\pi)\) and the angular bridge \(\psi\approx0.1503378808\) together constitute the tan-arctan-τ family. These quantities are transcendental; they are not rational multiples of \(\pi\). Their triangular content is already latent: the phase-slip identity \(\sin(\Delta\phi_{\rm torque})=\cos\psi\) and the orthogonal complement \(\Delta\phi_{\rm torque}-\psi=\pi/2\) encode the elementary 60-degree relations that will later crystallize into equilateral triangles. The family therefore carries triangular symmetry at the continuous level before any discrete orbit is formed.

When the bridge \(\psi\) is iterated seven times, the practical step  
\[
7\psi=\frac{\pi}{3}+\delta,\qquad\delta\approx0.0051676\ \text{rad}
\]  
appears. The residual \(\delta\) is small enough that six successive applications of this step partition the circle into six nearly equal sectors. The vertices of the resulting figure form a practical hexagon; the diagonals that join every second vertex form the two interlocking equilateral triangles of the hexagram. This is the completed radial fold—the geometric realization of the triangular symmetry that was latent in the continuous family. Because the residual is inherited uniformly by every sector, the practical hexagon is a uniformly expanded and slightly rotated copy of the ideal lattice. The six slices close after one full turn (modulo the small overshoot \(6\delta\)), preserving the radial-fold symmetry intact.

The completed radial fold does not remain a static ornament. Its six radial directions continue to act as a generator on the circle \(S^1\). Further multiples of the same practical step \(7\psi\) produce a dense sequence of points. Because \(7\psi/2\pi\) is irrational, the orbit is non-periodic: it never repeats exactly, yet it returns arbitrarily close to every previously visited direction. The residual six-sliced hexagon therefore propagates itself across a dense, non-periodic orbit on \(S^1\). Every new point still carries the same angular memory of the original six-fold partition; the radial-fold symmetry is transported along the dense trajectory without being erased by the density.

When this dense orbit is organized into concentric radial shells, the geometry of the outer hypercomplex projection appears. Ten, fifteen, and thirty points drawn from the same generator fill successive rings whose angular spacings approximate the Petrie projection of the \(E_8\) root system. The residual six-sliced hexagon supplies the local hexagonal coordination; the dense non-periodic filling supplies the progressive densification. The result is a nested sequence of Petrie rings whose innermost skeleton is the Order-6 starred polyhedra and whose outermost layers approach the 240-root geometry of \(E_8\). Thus the completed radial fold, once set in motion on the circle, propagates outward into the full hierarchy of nested shells.

Throughout the entire cascade the residual \(\delta\) remains the geometric signature of origin. It records that the six-fold structure was not imposed by an exact rational multiple of \(\pi\) but emerged from the continuous tan-arctan-τ family. Because that family is already the angular backbone of the Cosmological Clock—locked to the discrete tick lattice \(\varepsilon=10^{-9}\) and to the scale factor \(a(N)=\mathrm{e}^{\varepsilon N}\)—the residual, the radial fold, the dense orbit, and the \(E_8\) Petrie rings are all comoving. They advance with cosmic expansion, phase-locked at every epoch to the same microscopic separator that generated them.

In this manner seven consecutive comoving multiples of a continuous angular separator convert latent triangular symmetry into a completed radial fold; the residual six-sliced hexagon carried by that fold propagates across a dense non-periodic orbit on \(S^1\); and the same orbit, radially densified, generates the concentric Petrie rings that realize the outer \(E_8\) shell of the Decompositional Hierarchy.
